\section{Limitations}
\label{sec:limitations}

\begin{itemize}
  \item The monthly test block has 12 points.
        The weekly test block has 13 points.
        A model confidence set on those blocks is wide.
  \item Each grain has one test window.
        The study does not use a rolling origin.
        One window keeps the protocol the same for every model and every series.
        A rolling origin would answer a question about stability, not the question in this draft.
  \item Grain B keeps holidays and has no holiday indicator.
  \item E1 did not run.
        The TEM vendor tree is absent.
        The tables do not contain an E1 score.
  \item sMAPE on a series with many zero days is hard to read as an operational error.
        The rank still uses every day in grain A, including zeros.
  \item The skill score uses the Murphy ratio on sMAPE~\citep{murphy1988skill}.
        It compares a model only with the seasonal naive of the same grain.
  \item A low sMAPE can still miss the peaks.
        On the monthly total, boosting leads and the month-to-month order is weak.
        Read the scatter next to the rank.
  \item The Diebold-Mariano column compares the leader with the seasonal naive only.
        Pairwise comparisons of the other models are the Wilcoxon signed-rank tests in Table~\ref{tab:wilcoxon-pairs}, with a Benjamini--Hochberg adjustment inside the cell~\citep{benjamini1995fdr}.
        Those pairs are not a second rank.
  \item The HLN test in this paper uses horizon 1.
        Some models forecast the whole test block from one origin.
        Their errors are dependent.
        The p-value is a paired comparison of the stored losses, not a multi-step correction with \(h\) equal to the test length.
        That correction is not positive when \(h\) equals the test length.
  \item The Wilcoxon signed-rank test assumes independent differentials that are symmetric about the median.
        Forecast errors can be autocorrelated.
        The one-step HLN factor does not remove that dependence.
        With 12 or 13 test points the rank test has low power.
        The rank test is a robustness check.
        It does not replace the Diebold-Mariano test or the model confidence set.
\end{itemize}
